% Pursuing Stacks (cotes 134-2 to 134-8): archive cote 134-4, site page 81 (PDF page 82), typescript page 169.
% typescript. Draft compared: ps3.tex:4146-4216.
% AI-transcribed by Claude Opus 5.5 (claude-opus-5-5, adaptive thinking, effort high), Message Batches API, batch ps2 (msgbatch_01NVwr1Yuo8qJf5rykeGBXPK), 2026-11-12; not yet checked by eye.
\scanpage{134-4/81}
\mspage{169}

a) $W = u^{-1}(W')$ ( hence the functor $\bar{u} : W^{-1}M \to W'^{-1}M$ exists ) ,

b) the functor $\bar{u}$ is an equivalence .

We \unsure{de|do} not assume beforehand that $(M,W)$, $(M'\xout{p},W')$ are modelizers
(with respect to a given basic \unsure{localizer|hocalizer} $\underline{W} \subset F\unsure{l}((Cat))$), but in the ca-%
ses I have in mind, we'll know beforehand at least one of the pairs to
be a modelizer, and it will follow the other is one too.

In all what follows, a basic modelizer $\underline{W}$ is given once and for all
(\xout{dxfx} see definition on page 165). For the two main results\inserted{below} \deleted{on} the more
general test functors, we'll have to assume $\underline{W}$ satisfies the fibration
axiom L 5) (page 164). We are going to work with a fixed small catego-%
ry A, \xout{and} without any other preliminary assumptions upon A, all assump-%
tions that may be needed later will be stated in due course. Recall that
A can be considered as an object of (Cat), and we'll say A is $\underline{W}$-aspheric
if $A \to e$ is in $\underline{W}$ , which implies that $A \to e$ is even $\underline{W}$-aspheric i.e.
$A\mathrm{x}C \to C$ is in $\underline{W}$ for any C in (Cat). More generally, if F is in $A^{\wedge}$,
we'll call F $\underline{W}$-aspheric if the category $A_{/F}$ is $\underline{W}$-aspheric. Thus to
say A is $\underline{W}$-aspheric just means that the final object $e_{A^{\wedge}}$ of $A^{\wedge}$is $\underline{W}$-asphe-%
ric.

We'll constantly be using the canonical functor
\[ \hand{(1)} \qquad i_A : A^{\wedge} \longrightarrow (Cat) \quad , \qquad F \longmapsto A_{/F} \ , \]
and it's \xout{left} right adjoint
\[ \hand{(2)} \qquad j_A = i_A^{*} : (Cat) \longrightarrow A^{\wedge} \quad , \quad C \to (a \mapsto Hom(A_{/a} \ , C)). \]
The category $A^{\wedge}$ will always be viewed as endowed with the saturated set
of maps
\[ \hand{(3)} \qquad \underline{W}_A = i_A^{-1}(\underline{W}) \ . \]
This gives rise to the notions of $\underline{W}_A$-aspheric map in $A^{\wedge}$ , and of an ob-%
ject \xout{of} F of $A^{\wedge}$ being $\underline{W}_A$-aspheric over $e_{A^{\wedge}}$ namely $F \to e_{A^{\wedge}}$ being $\underline{W}_A$ -
aspheric, \xout{orxequivalently} which means that $F\mathrm{x}G \to G$ is in $\underline{W}_A$ for any
G in $A^{\wedge}$, which by definition\inserted{of $\underline{W}_{\unsure{A^{\wedge}|A}}$} means that for any G the map
\[ A_{/F\mathrm{x}G} \longrightarrow A_{/G} \]
in (Cat) is in $\underline{W}$ . Using the localization axiom on $\underline{W}$ , \xout{this} one sees
that it is enough to check this for G an object a of A , in which case
$A_{/G} = A_{/a}$ has a final object and hence is $\underline{W}$-aspheric by L 3), and the
condition amounts to $A_{/F\mathrm{x}a}$ being $\underline{W}$-aspheric, i.e. (with the termino-%
logy introduced above) that Fxa is $\underline{W}$-aspheric. Thus an object F of $A^{\wedge}$

\margin{81}
